% BHU PYQ -- Physics: Electromagnetic Theory & Relativistic Mechanics (2025-26 NEP)
% Paper Code: PHYMN41 / PHYMJ41 (Sem IV Mid-Semester Examination, Batch P2)
\documentclass[11pt,a4paper]{article}
\usepackage[a4paper,top=2.5cm,bottom=2.5cm,left=2.8cm,right=2.8cm,headheight=14pt]{geometry}
\usepackage{amsmath,amssymb}
\usepackage[T1]{fontenc}
\usepackage[utf8]{inputenc}
\usepackage{lmodern,microtype,fancyhdr,enumitem,hyperref,lastpage,parskip}

\hypersetup{colorlinks=true,urlcolor=black,linkcolor=black,pdftitle={PHYMN41: Electromagnetic Theory -- BHU Sem IV Mid-Term 2025-26 (NEP)}}

\pagestyle{fancy}\fancyhf{}
\renewcommand{\headrulewidth}{0.4pt}\renewcommand{\footrulewidth}{0.4pt}
\fancyhead[L]{\small \textit{Banaras Hindu University}}
\fancyhead[R]{\small \textit{PHYMN41: Electromagnetic Theory (Mid-Term 2025-26)}}
\fancyfoot[C]{\small Page \thepage\ of \pageref{LastPage}}

\newcommand{\pts}[1]{\hfill \textit{[#1]}}

\begin{document}

\begin{center}
  {\large\bfseries BANARAS HINDU UNIVERSITY}\\[2pt]
  {\normalsize B.Sc. (Hons.) Semester IV Mid-Semester Examination, 2025-26 (NEP)}\\[6pt]
  \rule{0.8\linewidth}{0.5pt}\\[6pt]
  {\large\bfseries DEPARTMENT OF PHYSICS}\\[4pt]
  {\large\bfseries Paper Code: PHYMN41 / PHYMJ41 : Electromagnetic Theory \& Relativistic Mechanics}\\[2pt]
  {\normalsize (Mid-Semester Examination --- Batch P2)}\\[4pt]
  {\normalsize Time: 1 Hour \hfill Maximum Marks: 20}
\end{center}
\rule{\linewidth}{0.4pt}
\smallskip

\noindent \textit{Instructions: Answer all questions from Section A and Section B as indicated. Figures to the right indicate full marks.}

\bigskip

\begin{center}
  {\large\bfseries SECTION A}
\end{center}

\bigskip

\noindent \textbf{Question 1.} Answer all parts:
\begin{enumerate}[label=(\alph*),itemsep=10pt,topsep=6pt]
  \item State divergence theorem. Give physical interpretation of the divergence theorem. \pts{4}
  
  \item Draw electric field lines produced by two equal and opposite point charges separated by some distance. \pts{2}
  
  \item Find the electric field a distance $z$ above the center of a circular loop of radius $r$ which carries a uniform line charge $\lambda$. \pts{4}
\end{enumerate}

\bigskip
\bigskip

\begin{center}
  {\large\bfseries SECTION B}
\end{center}

\bigskip

\begin{enumerate}[label=\textbf{\arabic*.},start=2,leftmargin=*,itemsep=14pt]
  \item A Michelson interferometer has arm length $L = 11\text{ cm}$, wavelength $\lambda = 500\text{ nm}$, and earth's velocity $v = 3 \times 10^4\text{ m/s}$. Calculate the expected fringe shift. Further, explain the negative (null) result obtained in Michelson-Morley experiment. \pts{3}

  \item Derive the Lorentz transformation equation starting from Einstein's postulates. \pts{4}
  \begin{center}\textbf{OR}\end{center}
  A rocket is chasing an enemy spaceship moving in the same direction. The enemy spaceship moves with a velocity of $0.6c$ relative to earth. The rocket moves with a velocity $0.8c$ relative to earth. Calculate the velocity of the rocket relative to the enemy spaceship using the relativistic velocity addition formula. \pts{4}

  \item Discuss the concept of zero rest mass of a photon and explain its implications in relativity. \pts{3}
\end{enumerate}

\bigskip
\begin{center}
  \textbf{***}
\end{center}

\end{document}
